% ModernSoftwarePractices.tex -- Modern Software Practices: A Condensed Guide (October 2026)
% An unofficial condensation of the DoW CIO library's software modernization, DevSecOps, open source and cATO documents.
% Sources: the documents listed in Appendix B (the user's copies, git-ignored under refs/; notes per document were kept
% outside the repository). Every claim carries the source's short name and PDF page.
% Build:  perl sim/msp.pl            (XeLaTeX via latexmk in a temporary directory -> docs/ModernSoftwarePractices.pdf)
%   or:   cd docs/src/msp && latexmk -xelatex ModernSoftwarePractices.tex
\documentclass[11pt,letterpaper,twoside,openany]{memoir}
\usepackage{fontspec}
\defaultfontfeatures[\rmfamily,\sffamily]{Ligatures=TeX}
\setmainfont{texgyrepagella}[Extension=.otf, UprightFont=*-regular, BoldFont=*-bold, ItalicFont=*-italic, BoldItalicFont=*-bolditalic]
\setsansfont{texgyreheros}[Extension=.otf, UprightFont=*-regular, BoldFont=*-bold, ItalicFont=*-italic, BoldItalicFont=*-bolditalic, Scale=MatchLowercase]
\setmonofont{DejaVuSansMono}[Extension=.ttf, BoldFont=*-Bold, ItalicFont=*-Oblique, BoldItalicFont=*-BoldOblique, Scale=0.84, Ligatures=NoCommon]
\newfontfamily\symfont{DejaVuSans}[Extension=.ttf, Scale=0.9]
\IfFontExistsTF{LibertinusSans-Regular.otf}
  {\newfontfamily\coverfont{LibertinusSans}[Extension=.otf, UprightFont=*-Regular, BoldFont=*-Bold]}
  {\newfontfamily\coverfont{TeX Gyre Heros}}
\usepackage{microtype}
\usepackage{xcolor,graphicx,array,booktabs,longtable,calc}
\usepackage{tikz}
\usetikzlibrary{positioning,arrows.meta,fit}
\usepackage[shortlabels]{enumitem}
\setlist{itemsep=1pt,topsep=4pt,parsep=0pt}
\usepackage[most]{tcolorbox}
\usepackage{multicol}
\usepackage{xurl}
\usepackage[hidelinks,bookmarksnumbered,pdfusetitle]{hyperref}
\usepackage{bookmark}
\graphicspath{{./}{../../img/}}
% ---- page layout: as the MBSE-DEVSECOPS Primer
\setlrmarginsandblock{1.15in}{1in}{*}
\setulmarginsandblock{1in}{1in}{*}
\setheadfoot{\onelineskip}{2\onelineskip}
\setheaderspaces{*}{1.5\onelineskip}{*}
\checkandfixthelayout
\setsecnumdepth{subsection}
\settocdepth{section}
\pagestyle{ruled}
\makeevenhead{ruled}{\small\thepage}{}{\small\sffamily\leftmark}
\makeoddhead{ruled}{\small\sffamily\rightmark}{}{\small\thepage}
\makeevenfoot{ruled}{}{\scriptsize\sffamily Unofficial condensation --- the original documents govern}{}
\makeoddfoot{ruled}{}{\scriptsize\sffamily Unofficial condensation --- the original documents govern}{}
\nouppercaseheads
\makechapterstyle{msp}{%
  \renewcommand*{\chapterheadstart}{\vspace*{1.2\onelineskip}}
  \renewcommand*{\chapnamefont}{\normalfont\sffamily\small\MakeUppercase}
  \renewcommand*{\chapnumfont}{\normalfont\sffamily\fontsize{36}{36}\selectfont\color{black!70}}
  \renewcommand*{\chaptitlefont}{\normalfont\sffamily\bfseries\huge\raggedright}
  \renewcommand*{\printchaptername}{\chapnamefont\chaptername}
  \renewcommand*{\chapternamenum}{\par\vspace{2pt}}
  \renewcommand*{\printchapternum}{\chapnumfont\thechapter}
  \renewcommand*{\afterchapternum}{\par\vspace{6pt}}
  \renewcommand*{\printchaptertitle}[1]{\chaptitlefont ##1\par\vspace{4pt}{\color{black!60}\rule{\textwidth}{0.6pt}}}
  \renewcommand*{\afterchaptertitle}{\par\vspace{1.2\onelineskip}}}
\chapterstyle{msp}
\setsecheadstyle{\normalfont\sffamily\bfseries\Large\raggedright}
\setsubsecheadstyle{\normalfont\sffamily\bfseries\large\raggedright}
\setsubsubsecheadstyle{\normalfont\sffamily\bfseries\raggedright}
\captionnamefont{\sffamily\bfseries\small}
\captiontitlefont{\sffamily\small}
\setlength{\LTcapwidth}{\linewidth}
\setlength{\parskip}{0.35\onelineskip}
\setlength{\parindent}{0pt}
\setlength{\emergencystretch}{3em}
\tolerance=2000
\hbadness=10000
\renewcommand{\arraystretch}{1.15}
% ---- conventions
\newcommand{\pg}[2]{{\small(#1~p.~#2)}}          % a source citation: short name and PDF page
\newcommand{\checkbox}{{\symfont\char"2610}\ }
\newcommand{\st}[1]{{\sffamily\scriptsize\bfseries #1}}   % a status tag
\tcbset{pbox/.style={enhanced,breakable,sharp corners,colback=white,colframe=black!65,boxrule=0.6pt,
  left=7pt,right=7pt,top=5pt,bottom=5pt,fonttitle=\sffamily\bfseries\small\raggedright,coltitle=black,colbacktitle=black!10,
  titlerule=0.4pt,toptitle=2pt,bottomtitle=2pt,before skip=10pt,after skip=10pt,parbox=false}}
\newtcolorbox{boxnote}[1]{pbox,title={#1}}
\newtcolorbox{boxteam}[1]{pbox,colback=black!3,frame hidden,borderline west={3pt}{0pt}{black!60},colbacktitle=black!3,titlerule=0pt,
  title={For a small MBSE team: #1}}
\newtcolorbox{boxwarning}[1]{pbox,colframe=black,boxrule=1.2pt,colbacktitle=black!18,title={#1}}
\newcolumntype{L}[1]{>{\raggedright\arraybackslash}p{#1}}
\setpnumwidth{2.8em}\setrmarg{3.8em}
\renewcommand{\cftchapterfont}{\sffamily\bfseries}
\title{Modern Software Practices}
\author{The agile kit}
\date{October 2026}
% ---- cover (tools/cover/cover.tex design: colour art, navy band with gold rules)
\definecolor{bandnavy}{HTML}{120E33}\definecolor{bandgold}{HTML}{D4A84B}\definecolor{bandgolddim}{HTML}{8F6B2A}\definecolor{bandcream}{HTML}{F6ECD6}
\newcommand{\CoverTitle}{Modern Software Practices}
\newcommand{\CoverLine}{A Condensed Guide}
\newcommand{\CoverSubtitle}{The DoW software modernization, DevSecOps, open source and cATO documents,\\ condensed for systems engineers and program teams}
\newcommand{\CoverDate}{October 2026}
\begin{document}
\frontmatter
\begin{titlingpage}
\thispagestyle{empty}
\begin{tikzpicture}[remember picture,overlay]
  \fill[white] (current page.south west) rectangle (current page.north east);
  \begin{scope}
    \clip (current page.south west) rectangle (current page.north east);
    \node[inner sep=0pt,anchor=center] at ([yshift=0.5in]current page.center) {\includegraphics[width=\paperwidth]{cover-art-color.jpg}};
  \end{scope}
  \fill[bandnavy,opacity=0.94] ([yshift=0.45in]current page.south west) rectangle ([yshift=2.95in]current page.south east);
  \draw[bandgold,line width=1.4pt] ([yshift=2.95in]current page.south west) -- ([yshift=2.95in]current page.south east);
  \draw[bandgold,line width=1.4pt] ([yshift=0.45in]current page.south west) -- ([yshift=0.45in]current page.south east);
  \draw[bandgolddim,line width=0.5pt] ([yshift=2.89in]current page.south west) -- ([yshift=2.89in]current page.south east);
  \draw[bandgolddim,line width=0.5pt] ([yshift=0.51in]current page.south west) -- ([yshift=0.51in]current page.south east);
  \node[bandcream,font=\coverfont\bfseries\fontsize{28}{33}\selectfont,align=center] (t) at ([yshift=2.30in]current page.south) {\CoverTitle};
  \node[bandgold,font=\coverfont,anchor=north] (s) at ([yshift=-0.10in]t.south) {\Large\scshape \CoverLine};
  \node[bandcream,font=\coverfont,anchor=north,align=center,text width=0.8\paperwidth] at ([yshift=-0.10in]s.south)
    {\large\hyphenpenalty=10000\exhyphenpenalty=10000 \CoverSubtitle\quad{\normalsize\color{bandgold}\textbullet\ \CoverDate}};
\end{tikzpicture}
\end{titlingpage}
\cleardoublepage
\thispagestyle{empty}
\vspace*{1.5in}
\begin{center}
{\sffamily\bfseries\Huge Modern Software Practices}\par\medskip
{\sffamily\Large A Condensed Guide}\par\bigskip
{\large The Department of War software modernization, DevSecOps, open source\\ and continuous authorization documents, as of October 2026}\par\vfill
{\small October 2026}
\end{center}
\clearpage
\thispagestyle{empty}
\vspace*{\fill}
\begin{boxwarning}{Read this first}
\textbf{This is an unofficial condensation. The originals govern.} It summarizes nineteen publicly listed documents of
the Department of War Chief Information Officer (DoW CIO) library, as linked on 3 October 2026 and read in the copies
listed, with their SHA-256 checksums, in Appendix~\ref{app:sources}. It adds no requirement of its own. Where it says
what a document means for a small systems-engineering team, the box is marked \emph{For a small MBSE team} and that
reading is this guide's, not the Department's.

\textbf{Names.} The Department of Defense now styles itself the Department of War. Documents published before the
change say ``DoD''; the library now lists several of them under ``DoW'' titles. This guide uses each document's own
name and the library's title for it.

\textbf{Citations.} ``\pg{Fundamentals}{12}'' means page 12 of the PDF file, counting the cover as page~1, which is
not always the number printed on the page. The short names are listed in Table~\ref{tab:status}.

\textbf{Status words.} \emph{Policy} binds (an instruction or a signed memorandum). \emph{Guidance} explains and
recommends. A \emph{report} records findings. \emph{Draft} is marked pre-decisional by the document itself. Several
2021 documents print ``This document automatically expires 1-year from publication date unless revised''; they are
still listed in the library, and this guide flags them.

\textbf{One marking to respect.} The cATO Implementation Guide carries Distribution Statement~C on its cover although it
also bears an open-publication clearance stamp and is posted in the public library. This guide summarizes it at a high
level only. Ask your security office before redistributing that original.
\end{boxwarning}
\vspace*{\fill}
\clearpage
\tableofcontents*
\clearpage
\mainmatter

% =====================================================================================================================
\chapter{Bottom line up front}\label{ch:bluf}
\vspace{-1\onelineskip}
Ten things a program team should know before reading any of the nineteen documents.

{\small
\begin{enumerate}[leftmargin=1.6em,itemsep=2pt]
\item \textbf{There is now an instruction.} DoWI 8430.01, \emph{Accelerated Mission Software}, effective 8 September
  2026, applies to all software throughout its lifecycle in every DoW Component, acquisition and non-acquisition
  programs alike, regardless of dollar value \pg{DoWI 8430.01}{1, 4}. Everything else in this guide is older and
  either directs one topic (the memos) or explains (guidance and reports).
\item \textbf{The direction is the same in every document:} small cross-functional teams, CI/CD pipelines, automation,
  security built in, continuous feedback from users \pg{DoWI 8430.01}{13--15}, \pg{Fundamentals}{9--10}.
  ``Treat testing, security scanning, compliance checks, and deployment as continuous activities, not final gates''
  \pg{DoWI 8430.01}{14}.
\item \textbf{Pipelines are critical infrastructure.} Component heads ``operate secure, automated CI/CD pipelines'' and
  treat them as infrastructure that ``must be authorized, monitored, and protected'' \pg{DoWI 8430.01}{9}.
\item \textbf{Evidence is machine-readable.} All software on DoW networks is to be supported by machine-readable
  evidence of its composition, build integrity and application security; an accurate SBOM for all covered software;
  provenance and attestations from an automated build \pg{DoWI 8430.01}{18}.
\item \textbf{The vocabulary is the Fundamentals.} Ten lifecycle phases (Plan to Feedback), software factory,
  DevSecOps platform, CI/CD pipeline, Infrastructure as Code \pg{Fundamentals}{9--12}.
\item \textbf{The checklist is the Activities \& Tools Guidebook.} Its REQUIRED activities ``are not negotiable'';
  a reference design may add activities or raise PREFERRED to REQUIRED, never lower one \pg{Guidebook}{6, 10}.
\item \textbf{Only two reference designs are approved.} CNCF Kubernetes (2021) and CNCF Multi-Cluster Kubernetes
  (2022). The AWS and the Azure + GitHub designs are drafts \pg{RD-AWS}{1}, \pg{RD-Azure}{2}.
\item \textbf{cATO is a privilege on top of an ATO.} It needs an existing ATO in the RMF monitor step, continuous
  monitoring, active cyber defense and an approved reference design; the DoD CISO approves it \pg{cATO memo}{1--3},
  \pg{cATO criteria}{6, 8}.
\item \textbf{Measure the value stream; use DORA metrics inside the team.} DORA metrics ``are not appropriate to compare
  beyond the team context''; enterprise reporting focuses on value streams and mission outcomes \pg{DoWI 8430.01}{16}.
\item \textbf{Open source and AI have rules.} Adopt, buy, then create \pg{OSS memo}{3}; assess third-party and open
  source components on sustainment, source, dependencies, security, integrity and foreign influence \pg{DoWI 8430.01}{22};
  AI-generated code is ``unverified input'' and non-public information stays out of unapproved AI services
  \pg{DoWI 8430.01}{23}.
\end{enumerate}
}

\begin{boxteam}{the short version}
{\small Use an existing DevSecOps platform; keep the model and its pipeline in it; let the pipeline produce the
evidence (tests, scans, SBOMs, the verification matrix) on every change; report value-stream measures upward and keep
DORA numbers inside the team.}
\end{boxteam}

% =====================================================================================================================
\chapter{How the documents fit together}\label{ch:map}

The library groups the documents under five headings. Read top-down, they form a chain: a strategy sets goals, an
implementation plan assigns tasks, an instruction makes them policy, the fundamentals give the vocabulary, the
guidebook and playbook give the activities, the reference designs give tested architectures, and the cATO documents
say how a platform built that way is authorized continuously.

\begin{figure}[htbp]
\centering
\begin{tikzpicture}[font=\sffamily\footnotesize, node distance=5mm and 6mm,
  b/.style={draw, rounded corners=1pt, align=center, text width=31mm, minimum height=11mm, fill=white},
  g/.style={draw=black!40, dashed, rounded corners=2pt, inner sep=2.5mm},
  a/.style={-{Stealth[length=2mm]}, black!70}]
\node[b] (strat) {Software Modernization\\Strategy (2021/22)};
\node[b, right=of strat] (plan) {Implementation Plan\\FY25--26 (2025)};
\node[b, right=of plan] (swft) {SWFT memo (2025)\\+ RFI summary};
\node[b, below=8mm of plan, text width=40mm, fill=black!8] (dowi) {\textbf{DoWI 8430.01}\\Accelerated Mission Software\\(policy, Sept 2026)};
\node[b, below=8mm of dowi] (fund) {DevSecOps\\Fundamentals (2024)};
\node[b, left=of fund] (guide) {Activities \& Tools\\Guidebook + xlsx (2025)};
\node[b, right=of fund] (play) {Playbook (2021)\\State of DevSecOps (2025)};
\node[b, below=8mm of guide] (path) {Pathway to a\\Reference Design};
\node[b, below=8mm of fund] (rd) {Reference designs:\\K8s, Multi-Cluster\\(drafts: AWS, Azure)};
\node[b, below=8mm of play] (oss) {Open Source\\memo (2022)};
\node[b, below=8mm of rd, text width=60mm] (cato) {cATO memo (2022) $\rightarrow$ Evaluation Criteria (2024)\\$\rightarrow$ Implementation Guide (2024)};
\draw[a] (strat) -- (plan); \draw[a] (plan) -- (swft);
\draw[a] (plan) -- (dowi); \draw[a] (strat.south) |- (dowi.west); \draw[a] (swft.south) |- (dowi.east);
\draw[a] (dowi) -- (fund); \draw[a] (fund) -- (guide); \draw[a] (fund) -- (play);
\draw[a] (guide) -- (path); \draw[a] (path) -- (rd); \draw[a] (guide.south east) -- (rd.north west);
\draw[a] (rd) -- (cato); \draw[a] (oss.south) |- (cato.east);
\end{tikzpicture}
\caption{How the nineteen documents relate (this guide's reading). Arrows point from the document that sets
direction to the one that implements or depends on it.}
\label{fig:map}
\end{figure}

\needspace{10\baselineskip}
\begin{longtable}{@{}L{0.28\linewidth}L{0.16\linewidth}L{0.14\linewidth}L{0.32\linewidth}@{}}
\caption{The nineteen documents: short name used in this guide, date, status.}\label{tab:status}\\
\toprule \sffamily\bfseries Document (library title) & \sffamily\bfseries Short name & \sffamily\bfseries Date & \sffamily\bfseries Status\\ \midrule\endfirsthead
\toprule \sffamily\bfseries Document (library title) & \sffamily\bfseries Short name & \sffamily\bfseries Date & \sffamily\bfseries Status\\ \midrule\endhead
\bottomrule\endfoot
\multicolumn{4}{@{}l}{\sffamily\bfseries Modern Software Practices}\\
DoW Software Modernization Strategy & Strategy & Nov 2021; memo 1 Feb 2022 & Strategy, directed by DepSecDef memo\\
Software Modernization Implementation Plan FY25--26 & I-Plan & cleared Apr 2025 & Plan (tasks); some pages carry inconsistent footer markings\\
DoWI 8430.01, Accelerated Mission Software & DoWI 8430.01 & effective 8 Sep 2026 & \st{POLICY} instruction\\
Memo: Accelerating Secure Software & SWFT memo & 24 Apr 2025 & \st{POLICY} memo (directs SWFT)\\
Software Fast Track RFI Combined Summary & SWFT RFI & cleared Dec 2025 & Report of industry responses, not an endorsement\\
\multicolumn{4}{@{}l}{\sffamily\bfseries DevSecOps document set}\\
DoW Enterprise DevSecOps Fundamentals & Fundamentals & v2.5, Oct 2024 & Guidance (approved by the SSG)\\
DoW Enterprise DevSecOps Activities \& Tools Guidebook & Guidebook & v2.5, Apr 2025 & Guidance with a REQUIRED baseline\\
DevSecOps Fundamentals Guidebook: Activities \& Tools (xlsx) & Workbook & companion of v2.5 & The guidebook's tables as a workbook\\
DevSecOps Playbook & Playbook & v2.1, Sep 2021 & Guidance; \st{SELF-EXPIRY} printed on cover\\
The State of DevSecOps Report & State report & Mar 2025 & Report (findings)\\
\multicolumn{4}{@{}l}{\sffamily\bfseries DevSecOps reference designs}\\
Pathway to Reference Design & Pathway & Sep 2021 & Process description\\
Reference Design: CNCF Kubernetes & RD-K8s & v2.1, Sep 2021 & Approved design; \st{SELF-EXPIRY} printed\\
Reference Design: CNCF Multi-Cluster Kubernetes & RD-Multi & v1.0, Jul 2022 & Approved design\\
DRAFT: DOW AWS Managed Service & RD-AWS & v0.2, Sep 2021 & \st{DRAFT}; approval ``TBD''; \st{SELF-EXPIRY}\\
DRAFT: DOW Cloud IaC + GitHub + Azure & RD-Azure & v0.01, Sep 2022 & \st{DRAFT} (signed, marked pre-decisional)\\
\multicolumn{4}{@{}l}{\sffamily\bfseries Open source software}\\
Software Development and Open Source Software Memorandum & OSS memo & 24 Jan 2022 & \st{POLICY} memo\\
\multicolumn{4}{@{}l}{\sffamily\bfseries Continuous authority to operate}\\
Memo: Continuous Authorization to Operate (cATO) & cATO memo & 2 Feb 2022 & \st{POLICY} memo (DoD SISO)\\
cATO Evaluation Criteria (DevSecOps use case) & cATO criteria & 29 May 2024 & Criteria; ``living document'' on the RMF KS\\
DevSecOps Continuous Authorization Implementation Guide & cATO guide & v1.0, Mar 2024 & Guidance; cover marked Distribution~C\\
\end{longtable}

\begin{boxnote}{Version drift between the documents}
The documents cite each other by version, and the versions moved on. The CNCF Kubernetes design asserts compatibility
with the Strategy Guide and the Tools and Activities Guidebook \emph{version 2.1} \pg{RD-K8s}{8}; the current guidebook
is version 2.5 \pg{Guidebook}{1}. The cATO guide cites the Fundamentals version 2.1 and the guidebook version 2.2
\pg{cATO guide}{26}. When two documents disagree on a detail, the later document and the instruction win; for an
authorization, ask your AO which version they assess against.
\end{boxnote}

% =====================================================================================================================
\chapter{Modern software practices: strategy, plan, policy}\label{ch:msp}

\section{DoW Software Modernization Strategy}
\textbf{What it is.} The DoD Software Modernization Strategy, version 1.0 of November 2021, approved by the Deputy
Secretary of Defense in a memo of 1 February 2022 \pg{Strategy}{1--2}. It is a sub-strategy of the Digital Modernization
Strategy and replaces the DoD Cloud Strategy \pg{Strategy}{3}. Audience: Department leadership and everyone who plans
software work.

\textbf{Key points.}
\begin{itemize}
\item Vision: ``Deliver Resilient Software Capability at the Speed of Relevance''; resilience means high-quality and
  secure software \pg{Strategy}{5--6}.
\item Five unifying principles: Primacy of Security, Stability, and Quality at Speed; Cloud Smart/Data Smart; Enterprise
  First; No One Left Behind; More Than Code (contracting, intellectual property, transition to fielding) \pg{Strategy}{6}.
\item Software is adopted, bought or created, and ``software delivery is not a one-and-done activity'' \pg{Strategy}{7}.
\item The framework pairs technical enablers (enterprise cloud, design patterns, DevSecOps tooling, enterprise
  services) with process transformation (business operations, acquisition, cyber survivability, testing, workforce)
  \pg{Strategy}{7--9}. A software factory is ``a software assembly plant'' with multiple pipelines producing deployable
  artifacts ``with minimal human intervention'' \pg{Strategy}{8}.
\item ``Cybersecurity should be the driver and compliance an outcome'' \pg{Strategy}{9}.
\item Three goals: accelerate the enterprise cloud environment (including automated design patterns such as
  Infrastructure as Code, Compliance as Code and hardened containers) \pg{Strategy}{10--11}; establish a
  Department-wide software factory ecosystem (including continuous authorization, where cyber professionals work with
  developers \emph{and system engineers} to identify pipeline-generated evidence) \pg{Strategy}{11--12}; transform
  processes (policy, agile acquisition, software as data, workforce) \pg{Strategy}{12--13}.
\item The CIO and the Under Secretaries for A\&S and R\&E lead through a Software Modernization Senior Steering Group
  (SSG) \pg{Strategy}{1, 14}.
\end{itemize}
\textbf{Requires or recommends.} A strategy: it sets direction and tasks the SSG; it does not itself bind programs.

\section{Software Modernization Implementation Plan, FY25--26}
\textbf{What it is.} The task list that carries the strategy forward, cleared for publication in April 2025
\pg{I-Plan}{1}. Of the 41 tasks in the FY23--24 plan, 27 were accomplished, 12 carried over and 2 combined
\pg{I-Plan}{4}. Audience: those responsible for an organization's software strategy, developers and operators, and
leaders \pg{I-Plan}{5}.

\textbf{Tasks that matter to a program team} (the plan lists them under the strategy's three goals, not in priority
order \pg{I-Plan}{6}):
\begin{itemize}
\item a quick-track Software-as-a-Service ATO process \pg{I-Plan}{7};
\item scale modern software practices, including a DevSecOps platform and software factory catalog and DoD-wide
  software assurance capabilities \pg{I-Plan}{9};
\item developer tools (web.git.mil, digital engineering for software development, Infrastructure as Code for the cloud
  contract's capabilities and at IL6) \pg{I-Plan}{10};
\item cATO ``must become a standard business practice'': designate cATO platforms, publish analytics, train
  authorizing officials \pg{I-Plan}{11};
\item prepare software factories for AI: data provenance, AI traceability, risk assessment, AI/ML DevSecOps guidance
  \pg{I-Plan}{12};
\item secure software standards: the NIST Secure Software Development Framework (SSDF) is already integrated into the
  Activities and Tools guidebook; secure coding practices and metrics, secure container standards and an SBOM
  repository pilot are next \pg{I-Plan}{14};
\item modernize JCIDS for DevSecOps, ``oftentimes the root cause is unclear requirements'', with recommended metrics
  such as DORA metrics \pg{I-Plan}{14};
\item software modernization in embedded weapons systems (Weapons Ignite toolkit) \pg{I-Plan}{15}; more use of the
  Software Acquisition Pathway, whose adoption grew about 50\% a year \pg{I-Plan}{16}; a software cadre and eight
  software engineering work roles \pg{I-Plan}{13, 16--17}.
\end{itemize}

\section{DoWI 8430.01, Accelerated Mission Software}\label{sec:dowi}
\textbf{What it is.} The DoW instruction on software modernization and management, approved 31 August 2026 by the DoW
CIO and effective 8 September 2026 \pg{DoWI 8430.01}{1}. It applies to all software throughout its lifecycle except
research and prototype funding (budget activities 1--4) until that software transitions to acquisition or operations
\pg{DoWI 8430.01}{4}. \st{POLICY}

\subsection{Policy in one page}
Software is ``a strategic asset to be managed as an enduring capability throughout its lifecycle''
\pg{DoWI 8430.01}{4}. DoW software efforts:
\begin{enumerate}[(1)]
\item employ modern, secure, human-centric practices such as Agile and DevSecOps;
\item use ``measurable and testable requirements that are traceable to mission critical functions'', including
  cybersecurity, cyber resilience and zero trust suitability;
\item use component analysis, telemetry, observability and automated monitoring;
\item prefer existing software, open source, COTS and SaaS before developing anything new;
\item examine business processes before developing software \pg{DoWI 8430.01}{4--5}.
\end{enumerate}

\subsection{The seven tenets}
\begin{longtable}{@{}L{0.30\linewidth}L{0.64\linewidth}@{}}
\caption{Modern software tenets \pg{DoWI 8430.01}{13--15}.}\\
\toprule \sffamily\bfseries Tenet & \sffamily\bfseries What it asks\\ \midrule\endfirsthead
\bottomrule\endfoot
Embrace change through continuous feedback & Software is never done; CI/CD and frequent deployment of working software are ``the primary mechanism for refining requirements''.\\
Optimize for flow, not phases & Remove friction; small, cross-functional, mission-focused teams decide technical choices, workflows and tooling within architecture and security baselines.\\
Build security and quality in, not on & The whole team owns it; approved reference designs, enterprise platforms and provisionally authorized IaC ``make the secure path the easy path''.\\
Automate to the greatest extent possible & Testing, scanning, compliance checks and deployment are continuous, ``not final gates''; feedback on every code change.\\
Architect for evolvability and scale & Loosely coupled, highly cohesive, modular, well-defined APIs, containers; simplicity and testability.\\
Default to enterprise reuse & Source, documentation, APIs, data schemas and SBOMs discoverable with unlimited (else government purpose) rights; military or space software is CTI from inception.\\
Align governance with speed & Adapt acquisition, cyber and T\&E processes; delegate decisions to the lowest practical level; verify through automated evidence.\\
\end{longtable}

\subsection{What it requires of a program}
\begin{itemize}
\item \textbf{Pipelines.} Operate secure, automated CI/CD pipelines as critical infrastructure \pg{DoWI 8430.01}{9}.
\item \textbf{Assurance evidence.} Static and dynamic analysis, fuzz testing, penetration testing and secure-code
  reviews, with continuous authorization strategies from inception \pg{DoWI 8430.01}{10}.
\item \textbf{Supply chain.} SBOMs generated for custom software and captured for COTS; provenance and pedigree tracked
  \pg{DoWI 8430.01}{10}. A secure, automated, observable build that emits machine-readable provenance and attestations
  \pg{DoWI 8430.01}{18}. All third-party components, including transitive dependencies, in the SBOM
  \pg{DoWI 8430.01}{23}.
\item \textbf{Users.} Human-centered design, regular contact with representative users, usage and satisfaction
  telemetry, value stream metrics in recurring product reviews \pg{DoWI 8430.01}{10}.
\item \textbf{Test.} Integrated, continuous T\&E that unifies DT\&E and OT\&E; DT\&E follows the DevSecOps Activities
  and Tools Guidebook; adversarial testing covers pipelines, IaC scripts, source repositories and running
  applications; findings go into the backlog and are tracked to remediation; periodic independent government cyber
  DT\&E is still required \pg{DoWI 8430.01}{19--20}. Testers accept validated automated artifacts, scans, performance
  data and telemetry as evidence \pg{DoWI 8430.01}{9}.
\item \textbf{Security baseline.} Factories in the enterprise digital arsenal are authorized at least moderate
  confidentiality and integrity; cloud hosting at IL2 or higher; software and technical artifacts with a military or
  space application (source, build scripts, SBOMs, test vectors, design documents) are CTI at moderate/moderate
  \pg{DoWI 8430.01}{20}.
\item \textbf{Common security requirements} in eight areas: secure architecture, authentication and authorization,
  input validation, FIPS-validated cryptography, error handling and logging, software safety, operational resilience,
  continuous vulnerability management with analysis and dependency scanning in pipelines \pg{DoWI 8430.01}{20--21}.
\item \textbf{Measurement.} Value stream metrics across the full lifecycle; team metrics such as DORA are leading
  indicators ``not appropriate to compare beyond the team context''; contracts give the government access to the
  metrics data in contractor environments \pg{DoWI 8430.01}{15--16}.
\item \textbf{Reuse and release.} Report custom code (DITPR for unclassified); share it in a repository after removing
  PII and sensitive data ``from code, documentation, issues, and commit history''; military or space code is shared by
  InnerSource; other code may be released as open source under permissive, widely recognized licences
  \pg{DoWI 8430.01}{25--26}.
\item \textbf{Acquisition.} The Software Acquisition Pathway (DoDI 5000.87) is ``the preferred pathway for all DoW
  software acquisition''; commercial solutions openings and other transactions are the default solicitations
  \pg{DoWI 8430.01}{27}.
\end{itemize}

\subsection{AI in software development}
Developers stay fully accountable; security- or safety-critical AI-generated changes need human review and approval;
AI-generated code is ``unverified input''; non-public DoW information --- code, configuration scripts, infrastructure
definitions, schematics, documentation --- may not go into generative AI unless the service resides on DoW information
systems and is approved; AI code gets the same review and testing as other code, including intellectual property and
licence checks; only approved AI services, which must not train on government data or prompts and must be auditable;
keep a record of models, versions and datasets in the evidence package ``analogous to the SBOM'' \pg{DoWI 8430.01}{23--25}.

\begin{boxteam}{DoWI 8430.01}
A model repository in GitLab with a pipeline is already most of tenets (a), (c) and (d). The instruction's phrase
``requirements that are traceable to mission critical functions'' is the requirement-to-verification trace a SysML v2
model keeps (Primer: \emph{The pipeline, job by job}). The AI rules apply to a local model assistant too: keep the model
on an approved system, review what it writes, and record which model wrote it (Primer: \emph{A local AI assistant on
the laptop}). Removing sensitive data ``from \dots commit history'' before sharing is a reason to keep it out of the
history in the first place.
\end{boxteam}

\section{Memo: Accelerating Secure Software}
\textbf{What it is.} A one-page memo of 24 April 2025 by the official performing the duties of the CIO, directing a
Software Fast-Track (SWFT) initiative \pg{SWFT memo}{1}. \st{POLICY}

It names the problem plainly: ``Lengthy, outdated cybersecurity authorization processes frustrate agile, continuous
delivery'', and the Department ``lacks visibility into the origins and security of software code''. SWFT is to define
clear cybersecurity and supply chain risk management requirements, rigorous software security verification, secure
information sharing, and government-led risk determinations to expedite authorizations; the framework and plan were
due within 90 days \pg{SWFT memo}{1}.

\section{Software Fast Track RFI Combined Summary}
\textbf{What it is.} An analysis of 420 industry responses to three requests for information on supply chain security
tools, external assessment, and automation and AI, cleared in December 2025 \pg{SWFT RFI}{1, 3}. It ``does not represent
an endorsement \dots nor dictate what a SWFT implementation plan will encompass'' \pg{SWFT RFI}{7}, and it says it was
partially assisted by generative AI with human review \pg{SWFT RFI}{1}. \st{REPORT}

\textbf{Findings worth knowing.} NIST SP 800-218 (SSDF) appeared in over 75\% of responses, OWASP material in 66\%,
NIST SP 800-53 in 60\%, NIST SP 800-161 in 46\% \pg{SWFT RFI}{7}; respondents asked the Department to define what a
valid SSDF attestation is \pg{SWFT RFI}{8}; 90\% would provide an SBOM \pg{SWFT RFI}{10}; risk artifacts mapped to the
RMF body of evidence (SAR, SSP, RAR, POA\&M) \pg{SWFT RFI}{11--12}; for SBOMs, standardized machine-readable formats
such as CycloneDX or SPDX, package URLs, licences, provenance and dependency relationships \pg{SWFT RFI}{26}.

% =====================================================================================================================
\chapter{The DevSecOps document set}\label{ch:dso}

\section{DoW Enterprise DevSecOps Fundamentals}
\textbf{What it is.} An ``educational compendium'' approved by the SSG on 16 October 2024 \pg{Fundamentals}{3, 7},
for novice and intermediate staff and for managers and leaders \pg{Fundamentals}{7}. \st{GUIDANCE}

\textbf{Definition.} DevSecOps is ``a combination of software engineering methodologies, practices, and tools that
unifies software development (Dev), security (Sec), and operations (Ops)'' with collaboration, automation and continuous
monitoring \pg{Fundamentals}{9}. DoWI 8430.01 uses the same definition \pg{DoWI 8430.01}{31}.

\textbf{Assumptions.} Agile is already in place; approved or provisionally authorized cloud is preferred, with
on-premises hardware-in-the-loop testing for weapons systems; ``Avoiding vendor lock-in without considering other types
of lock-in is ill-advised'' \pg{Fundamentals}{7--8}. A DevSecOps implementation ``does not require a specific
architecture, containers, or even explicit use of cloud computing'', though these are strongly recommended and some
reference designs mandate them \pg{Fundamentals}{12}.

\textbf{Implementation components.} The software supply chain (four links each 90\% secure give roughly 65\%; raising
each to 95\% gives 81\%) \pg{Fundamentals}{12--13}; the software factory (infrastructure, digital platform and
application capabilities) \pg{Fundamentals}{13--14}; the DevSecOps platform \pg{Fundamentals}{17--18}; CI/CD pipelines
\pg{Fundamentals}{18--20}; Infrastructure as Code, which can shorten deployment of networking, identity and security
policies ``from the standard 30 weeks to hours'' \pg{Fundamentals}{20}. A merge request's peer review is ``the software
code equivalent of two-person integrity'' \pg{Fundamentals}{15}; ``Released is never equivalent to Deployed!''
\pg{Fundamentals}{16}.

\textbf{Metrics.} DORA's four (deployment frequency, change lead time, change failure rate, failed deployment recovery
time), Google's four golden signals, the OUSD(A\&S) ``DoD4'' (speed and cadence, stability and reliability,
value/ROI, cyber resilience), OKRs and flow metrics (cycle time, throughput, work in progress)
\pg{Fundamentals}{33--35}.

\textbf{Culture.} Small incremental changes, security from the start, value open source, engage users early, automate
repeated manual work, ``Fail fast, learn fast, but don't fail twice for the same reason'', documentation close to the
code \pg{Fundamentals}{30--31}. Adoption is ``a marathon, not a sprint'' \pg{Fundamentals}{37}.

\section{Activities \& Tools Guidebook and workbook}\label{sec:guidebook}
\textbf{What it is.} Version 2.5 of April 2025 and its companion spreadsheet, ``a single package''
\pg{Guidebook}{1, 6}. The spreadsheet is also a template for submitting new reference designs and for ``building a
continuous ATO package that demonstrates process and responsibility mapping'' \pg{Guidebook}{6}. \st{GUIDANCE} with a
REQUIRED baseline.

\textbf{The baseline column.} REQUIRED activities ``are not negotiable'' and must be performed within the software
factory as part of the minimum viable product release; PREFERRED are aspirational; AS REQUIRED are at the mission owner
and AO's discretion \pg{Guidebook}{6, 10}. Reference designs cannot lower REQUIRED to PREFERRED; a contract, requirement
or reference design may make an activity AS REQUIRED or not applicable \pg{Guidebook}{10}. Each activity is mapped to
the SSDF practice groups PO, PS, PW and RV \pg{Guidebook}{11}. DevSecOps ``is most frequently adopted using Scrum 1 and 2
week Sprints'' \pg{Guidebook}{7}.

\textbf{Configuration management.} ``Without configuration management discipline, DevSecOps practices will not reach
their full potential'': identification, control (CMDB, source and artifact repositories), and verification and audit as
pipeline control gates \pg{Guidebook}{15--16}.

\begin{longtable}{@{}L{0.12\linewidth}L{0.30\linewidth}L{0.52\linewidth}@{}}
\caption{The ten lifecycle phases \pg{Fundamentals}{10} and their REQUIRED activities (Guidebook Tables 4--13,
PDF pp.~17--52; workbook phase tabs). Mission-based cyber risk assessment and test audit are REQUIRED in nearly every
phase and are not repeated.}\label{tab:phases}\\
\toprule \sffamily\bfseries Phase & \sffamily\bfseries Purpose & \sffamily\bfseries REQUIRED activities\\ \midrule\endfirsthead
\toprule \sffamily\bfseries Phase & \sffamily\bfseries Purpose & \sffamily\bfseries REQUIRED activities\\ \midrule\endhead
\bottomrule\endfoot
Plan & Requirements and objectives, focused on the next release & Change management planning; configuration identification; CM planning; DevSecOps process design; documentation version control; IaC deployment; project/release planning; team onboarding planning; risk management; software requirement analysis; system design; test deployment; test plan. (Threat modeling and design review are PREFERRED.)\\
Develop & Create the elements & Code commit; code commit logging; code commit scan; database component test; documentation; functional test; security code development; static analysis; static code scan before commit; test development; unit test. (Code review is PREFERRED.)\\
Build & Compile or integrate & API security tests; build; build configuration control and audit; component, functional, integration, regression and software integration tests; dependency vulnerability checking; release packaging; SAST; store artifacts.\\
Test & Verify against requirements & API security tests; compliance scan; manual security test; regression test; SBOM software composition analysis; service security test; software factory risk continuous monitoring; software integration test; SAST; system test. (DAST and performance test are PREFERRED.)\\
Release & Package and document & Operational readiness test; release go/no-go decision; SBOM SCA; software factory risk continuous monitoring. (User story review and demonstration is PREFERRED.)\\
Deliver & Make available to operations & Configuration integration testing; deliver released artifacts; delivery results review; operations team acceptance; SBOM SCA; factory risk monitoring.\\
Deploy & Install and configure & Compliance tests; interoperability tests; mission-oriented developmental tests; performance tests; post-deployment security scan; SBOM SCA; factory risk monitoring; user evaluation/feedback. (Formal qualification, OT\&E and others AS REQUIRED.)\\
Operate & Use in operations & Business operations; cyber OT\&E; logging; SBOM SCA; factory risk monitoring; user evaluation/feedback.\\
Monitor & Observe and measure & Asset inventory including SBOMs; compliance monitoring (COTS; resources and services); log analysis and auditing; system security monitoring; user evaluation/feedback.\\
Feedback & Changes for the next iteration & Revise product backlog; user evaluation/feedback.\\
\end{longtable}

\begin{boxteam}{the guidebook}
Several Plan-phase REQUIRED activities are systems engineering by another name: configuration identification, CM
planning, risk management, software requirement analysis, system design, test plan. A model-based team that keeps its
requirements, design and verification in a configuration-managed model with a pipeline already performs them; map them
in the workbook's tailoring columns rather than writing new documents. The Tools tab's ``requirements database'' asks to
``trace requirements to their source'' and ``trace features or design components to the requirements'' --- the model's
trace (Primer: \emph{DevSecOps in systems-engineering words}).
\end{boxteam}

\section{DevSecOps Playbook}
\textbf{What it is.} Eleven one-page plays with checklists, version 2.1 of September 2021, printed with a one-year
self-expiry \pg{Playbook}{1}. \st{GUIDANCE} \st{SELF-EXPIRY}

The plays: adopt a DevSecOps culture (``from `I' to `we'{}'') \pg{Playbook}{4}; adopt Infrastructure as Code (vendor
lock-in versus product lock-in: cloud-agnostic tools such as Ansible and Terraform ``avoid vendor lock-in but create product lock-in''; GitOps =
IaC + merge requests + CI/CD) \pg{Playbook}{5}; containerized microservices on CNCF Certified Kubernetes with Iron Bank
images, a sidecar security stack and a service mesh \pg{Playbook}{6}; a capability model, not a maturity model (DORA's
four key metrics) \pg{Playbook}{7}; 24 key capabilities \pg{Playbook}{8}; establish a software factory --- adopt an
existing one before building a new, expensive, multi-tenant one \pg{Playbook}{9}; define a meaningful pipeline (start
small; ``work closely with the AO to understand precisely what each control gate must validate'') \pg{Playbook}{10};
agile acquisition (``Do not lock technical requirements into legal contracts'') \pg{Playbook}{11}; cyber resilience
(``Do not use Fast Track Authority to Operate for software produced by a DevSecOps software factory CI/CD pipeline'';
help the AO move to a near real-time dashboard) \pg{Playbook}{12}; shift test and evaluation left (``Establish the plan
to automate the collection of test data metrics in the first sprint'') \pg{Playbook}{13}; industry practices
\pg{Playbook}{14}.

\section{The State of DevSecOps}
\textbf{What it is.} The first DoD CIO report on the state of DevSecOps, March 2025, based on interviews with more than
75 leaders and practitioners in 19 software organizations \pg{State report}{1, 4}. \st{REPORT}

\textbf{Findings.} Over 50 software factories deliver code to production \pg{State report}{3}; about 78 acquisition
programs had adopted the Software Acquisition Pathway, three quarters of them delivering software in under six months
\pg{State report}{10}; ``Software Factories Are Our Digital Arsenal'' \pg{State report}{4}; DevSecOps ``must be
introduced with leadership commitment'' \pg{State report}{3}. On cATO: a system needs an existing ATO and the RMF
monitoring stage; authorizing officials are often ``biased to caution'' rather than ``biased to action''
\pg{State report}{27--28}; reciprocity is now the default stance ``except when cybersecurity risk is too great''
\pg{State report}{29}. Policy should be S-P-E-E-E-D: Socialized, Portfolio-wide, Executable, Expertise-aware, Evolving,
Discoverable; ``Culture interprets policy'' \pg{State report}{32--33}.

\textbf{Its playbook for measurement} \pg{State report}{39--41}: do not average metrics across unlike systems; ask did we
build the right thing, did we build it right, did we get it to the right people, quickly and adaptably; and four rules ---
``The best data is the data used every day'', ``Manage to mission value, not metrics'', ``Don't rely on a single
metric'', ``Data can be aggregated, but metrics can't''.

% =====================================================================================================================
\chapter{DevSecOps reference designs}\label{ch:rd}

\section{Pathway to a Reference Design}
Eight steps from community proposal to approved design: proposal, evaluation and differentiation, define interconnects,
publish draft, execute pilot, lessons learned, SSG briefing, approved reference design \pg{Pathway}{2}. A draft must
build on the guidebook's REQUIRED items and ``must be piloted in an operationally relevant environment'' before approval;
approved designs are living documents ``revised no less than once per year'' with an expiration banner
\pg{Pathway}{3}.

\section{The four designs side by side}
{\footnotesize
\begin{longtable}{@{}L{0.15\linewidth}L{0.18\linewidth}L{0.18\linewidth}L{0.18\linewidth}L{0.18\linewidth}@{}}
\caption{Reference design comparison (each design's own statements).}\label{tab:rd}\\
\toprule & \sffamily\bfseries CNCF Kubernetes & \sffamily\bfseries Multi-Cluster K8s & \sffamily\bfseries AWS Managed (draft) & \sffamily\bfseries Azure + GitHub (draft)\\ \midrule\endfirsthead
\toprule & \sffamily\bfseries CNCF Kubernetes & \sffamily\bfseries Multi-Cluster K8s & \sffamily\bfseries AWS Managed (draft) & \sffamily\bfseries Azure + GitHub (draft)\\ \midrule\endhead
\bottomrule\endfoot
Status & Approved, v2.1, Sep 2021; self-expiry printed \pg{RD-K8s}{1, 3} & Approved, v1.0, Jul 2022 \pg{RD-Multi}{1, 3} & Pre-decisional draft v0.2; approval TBD \pg{RD-AWS}{1, 3} & Pre-decisional draft v0.01 (signed Oct 2022) \pg{RD-Azure}{1--2}\\
Stance & Cloud-agnostic; container is the unit of deployment; K8s in production \pg{RD-K8s}{7} & Many small clusters, shared services, forward-deployable management \pg{RD-Multi}{7, 10} & ``not vendor or product-agnostic''; deep AWS managed services \pg{RD-AWS}{9--10} & CSP-native SaaS + IaC; ``does not cover containers'' \pg{RD-Azure}{7, 9}\\
Platform & CNCF Certified K8s, Iron Bank hardened containers \pg{RD-K8s}{11, 18} & Conformant K8s with Cluster API lifecycle \pg{RD-Multi}{9} & EKS, ECR, App Mesh \pg{RD-AWS}{12--13} & DoD Cloud IaC on Azure Government; GitHub \pg{RD-Azure}{8--9}\\
North--south & Cloud Native Access Point \pg{RD-K8s}{11} & Delegated authentication within an accreditation boundary \pg{RD-Multi}{8} & CNAP \pg{RD-AWS}{12} & CNAP via Azure AD, Firewall, App Gateway/WAF or Front Door \pg{RD-Azure}{19}\\
East--west & Service mesh; sidecar security stack mandatory \pg{RD-K8s}{11} & mTLS between and within clusters \pg{RD-Multi}{10, 17} & App Mesh mTLS \pg{RD-AWS}{15} & Private endpoints, firewall hub \pg{RD-Azure}{52}\\
CI/CD & Orchestrator with stage entrance/exit control gates \pg{RD-K8s}{18--19} & Continuous reconciliation from a central source of truth \pg{RD-Multi}{12--14} & Pipeline with the main branch; CDK/CloudFormation \pg{RD-AWS}{11} & GitHub Actions; GitHub-hosted or self-hosted runners \pg{RD-Azure}{23, 30--33}\\
Monitoring & SIEM/SOAR; CSP services ``both/and'' \pg{RD-K8s}{30--31} & Production-ready K8s even in development \pg{RD-Multi}{10, 19} & CloudEvents, EventBridge, Security Hub \pg{RD-AWS}{13} & Monitor, Defender for Cloud, Policy, Sentinel \pg{RD-Azure}{14--15}\\
Notable & Two different container scanners REQUIRED \pg{RD-K8s}{25} & ``No in-place patching'' \pg{RD-Multi}{10} & No SSH or RDP; HTTPS only \pg{RD-AWS}{11} & ``code low, deploy high'' (IL2 dev, IL4/5 prod) \pg{RD-Azure}{8}\\
\end{longtable}}

\section{The Azure + GitHub draft in more detail}
Because the Azure draft is the closest to a cloud-hosted, SaaS-tooled team, its specifics are worth knowing even though
it is a draft and names GitHub rather than any other tool.
\begin{itemize}
\item Three design elements: Infrastructure as Code builds the cloud environment; a CSP-managed development suite as a
  service; applications from CSP-native services secured by policy \pg{RD-Azure}{8--9}.
\item The DoD Cloud IaC baseline has a DISA provisional ATO; inheritance from the PaaS and IaaS layers ``accounts for
  approximately 40\% of all controls''; systems ``must still be assessed for a full ATO'' \pg{RD-Azure}{7, 10, 19}. The
  baselines are updated monthly and are Distribution~C (US Government and contractors) \pg{RD-Azure}{20}.
\item Prerequisites: an empty Azure AD tenant, an Enterprise Agreement with five named subscriptions (Operations,
  Development, Test, Production, Forensics); pay-as-you-go subscriptions are incompatible \pg{RD-Azure}{16}.
\item Runners: GitHub-hosted runners in Azure Commercial for workloads up to IL5, or self-hosted runners inside the
  Azure Government environment, which poll outward over authenticated HTTPS so artifacts never leave the secured
  environment \pg{RD-Azure}{30--33}. Use OpenID Connect federation instead of long-lived cloud credentials in the CI
  secret store \pg{RD-Azure}{34--35}.
\item Pipeline minimums: linters, security checks, code coverage, functional checks; a periodic (weekly or daily) build
  so scanning happens even without changes \pg{RD-Azure}{29--30}. IL6 and air-gapped clouds need a self-hosted server
  product \pg{RD-Azure}{70}.
\end{itemize}

\begin{boxteam}{GitLab Dedicated on Azure with AKS runners}
No reference design in the library names GitLab. The cATO guide lists ``an existing commercial integrated DevSecOps
tool'' among its options for a platform \pg{cATO guide}{15}, and the Azure draft's patterns carry over by analogy:
runners inside your own Azure subscription that poll outward (a GitLab runner with the Kubernetes executor on AKS does
this), short-lived federated credentials instead of stored secrets, private repositories at IL4 and above, a scheduled
pipeline that scans even when nothing changed. Which reference design your platform claims conformance to is a
question for your platform team and AO, because cATO requires one \pg{cATO criteria}{14} (Primer: \emph{Runners: where
the pipeline runs}). Terraform is named in the Playbook as cloud-agnostic IaC that trades vendor lock-in for product
lock-in \pg{Playbook}{5}.
\end{boxteam}

% =====================================================================================================================
\chapter{Open source software}\label{ch:oss}

\textbf{What it is.} The DoD CIO memorandum of 24 January 2022, which supersedes the 2009 open source memo
\pg{OSS memo}{1--2}. \st{POLICY} DoWI 8430.01 adds the 2026 rules (sections 3.5.i and 3.7) \pg{DoWI 8430.01}{21--26}.

Two concerns frame it: externally maintained code is a path for malicious code, so open source needs the same supply
chain risk management as any product; and imprudent sharing of DoD code helps adversaries, which a modular open systems
approach manages by separating the critical parts \pg{OSS memo}{1}.

\begin{longtable}{@{}L{0.62\linewidth}L{0.32\linewidth}@{}}
\caption{Open source rules.}\label{tab:oss}\\
\toprule \sffamily\bfseries Rule & \sffamily\bfseries Source\\ \midrule\endfirsthead
\toprule \sffamily\bfseries Rule & \sffamily\bfseries Source\\ \midrule\endhead
\bottomrule\endfoot
Adopt, Buy, Create: existing government or open source first, then proprietary, new software last. & \pg{OSS memo}{3}; \pg{DoWI 8430.01}{10}\\
Open source is commercial computer software and gets equal consideration; refusing it only because it is open source may restrict competition. & \pg{OSS memo}{3}\\
Program managers own the suitability of off-the-shelf components, including open source. & \pg{OSS memo}{3}\\
Assess: long-term support (project health), trusted sources, dependencies, component security, integrity (signatures), foreign influence. & \pg{OSS memo}{4--5}; \pg{DoWI 8430.01}{22}\\
For non-trivial open source, also governance and control model, community health and security practices, licence and legal risk. & \pg{DoWI 8430.01}{22--23}\\
Respect the licence: most open source is copyrighted; keep authorship and licence notices. & \pg{OSS memo}{5}\\
Do not unduly restrict developers' access to software repositories. & \pg{OSS memo}{5}\\
Plan sustainment of open source components in the product support strategy. & \pg{OSS memo}{6}\\
Take delivery of source code and everything needed to build, test and deploy it. & \pg{OSS memo}{6}\\
Non-NSS software is open by default unless it is critical technology, the government lacks rights, or export rules forbid release. & \pg{OSS memo}{7}\\
Distribution B--F software is CTI and may not be released as open source; software that is not is not CTI. & \pg{OSS memo}{7}\\
Acceptable licences: Apache-2.0, BSD, GPL, LGPL, MIT; no custom licences. (DoWI: permissive, widely recognized licences approved by a Military Department or the CIO.) & \pg{OSS memo}{7--8}; \pg{DoWI 8430.01}{26}\\
Contribute upstream by pull/merge request; avoid DoD-specific forks; employees may contribute as official duty after consulting their supervisor. & \pg{OSS memo}{8}\\
All third-party components and transitive dependencies go in the SBOM. & \pg{DoWI 8430.01}{23}\\
\end{longtable}

\begin{boxteam}{open source tools in a model pipeline}
A team's toolchain (git, a scripting language, a SysML v2 grammar, a diagram tool) is open source under these rules:
record each in the SBOM, assess its project health once, and keep the tool versions in the baseline. A kit written by
the government team for a civilian, non-critical purpose is a candidate for release; anything with a military
application is CTI from its inception \pg{DoWI 8430.01}{14}.
\end{boxteam}

% =====================================================================================================================
\chapter{Continuous authorization to operate}\label{ch:cato}

\section{The memo}
Signed 2 February 2022 by the DoD Senior Information Security Officer \pg{cATO memo}{3}. \st{POLICY} RMF practice focused
on obtaining ATOs and fell short on continuous monitoring \pg{cATO memo}{1}. To achieve cATO the AO must demonstrate three
competencies: on-going visibility with robust continuous monitoring of RMF controls; active cyber defense in real time;
and ``the adoption and use of an approved DevSecOps reference design'' \pg{cATO memo}{1}. ``Simply conducting scans and
patching does not meet the threshold for active cyber defense'' \pg{cATO memo}{2}. The AO and component CISO present
the request and evidence to the DoD CISO; approved cATOs ``do not have an expiration date'' while the posture holds;
the cATO does not replace the underlying ATO and can be revoked without losing it \pg{cATO memo}{2--3}.

\section{The evaluation criteria and the implementation guide}
\textbf{Two use cases} \pg{cATO criteria}{5--6}: (1) inside the software factory boundary --- the factory's own
production environment, ``the main use case''; (2) outside it --- software deployed into another authorization
boundary (for example a weapon system) under a memorandum of understanding and interconnection security agreement.

\textbf{Prerequisites.} All RMF steps followed and the system in the RMF Monitor step \pg{cATO criteria}{6}; a current
ATO ``with no `High' or `Very High' unmitigated findings'' \pg{cATO criteria}{8}; a robust continuous monitoring
strategy with automated triggers on approved thresholds \pg{cATO criteria}{7}. Cloud hosting ``requires the deployment
of a Cloud Native Application Protection Platform (CNAPP)'' \pg{cATO criteria}{7}.

\textbf{Assessment.} Platform, process and people are assessed against the three competencies \pg{cATO criteria}{9};
evidence includes demonstrations, interviews and tracking systems \pg{cATO guide}{9--10}. The guide's practical advice:
use an existing platform before building one --- building new ``is the most time-consuming approach, and it should be
avoided if possible''; bring security into the team at the start; continuous monitoring in all environments, including
development and test; user stories and frequent demonstrations ``rather than creating comprehensive requirements''; no
waterfall and no heavy-weight agile process \pg{cATO guide}{15--16}.

\textbf{Metrics.} Mean time to patch; guardrail and control-gate trends sprint over sprint; feedback frequency;
mitigation effectiveness; security posture dashboards \pg{cATO guide}{14--15}, \pg{State report}{30}.

\section{Checklist}
Condensed from the evaluation criteria \pg{cATO criteria}{10--17}. The criteria mark the package documents; this list
is a reading aid, not the package.

\begin{longtable}{@{}L{0.04\linewidth}L{0.90\linewidth}@{}}
\caption{cATO evaluation checklist.}\label{tab:cato}\\
\toprule \multicolumn{2}{@{}l}{\sffamily\bfseries Continuous monitoring \pg{cATO criteria}{10--13}}\\ \midrule\endfirsthead
\bottomrule\endfoot
\checkbox & cATO risk management strategy with risk tolerances, insider and external threats (DoDI 8510.01)\\
\checkbox & System CONMON strategy: vulnerabilities tracked on all assets; implementation, effectiveness and impact measures; monitoring timelines; tied to audit and incident response\\
\checkbox & Authorization boundary diagram with data flows (including PII) and all external connections\\
\checkbox & Business rules: certified cyber staff, roles, a vulnerability POC, AO real-time access to test, scan and monitoring results\\
\checkbox & Automated monitoring: a working near-real-time dashboard, alerting, reporting to CMRS, the pipeline itself monitored\\
\checkbox & Authorization package (SAP, SAR, RAR, SSP, signed ATO, POA\&Ms), with a security assessment covering the platform and delivery pipeline (threat modeling, pen testing, attack surface, manual code review, SAST, DAST, IAST)\\
\checkbox & COOP/disaster recovery with test evidence; incident response plan with exercises\\
\checkbox & Continuous vulnerability management with timelines; audit log analysis; the cATO approval memo template\\
\midrule \multicolumn{2}{@{}l}{\sffamily\bfseries Active cyber defense \pg{cATO criteria}{13--14}}\\ \midrule
\checkbox & Certified CSSP with SLA, sensors, and training on DevSecOps\\
\checkbox & Third-party penetration test of development and operational environments within 90 days and annually\\
\checkbox & Ongoing security testing against real adversary techniques; budget for automated cyber testing\\
\midrule \multicolumn{2}{@{}l}{\sffamily\bfseries Platform \pg{cATO criteria}{14--15}}\\ \midrule
\checkbox & Name the approved reference design the platform implements\\
\checkbox & SBOM for the platform and an automated SBOM export for every product (SPDX, SWID or CycloneDX; the format ``has not yet been mandated''); an SBOM archive; how new CVEs are applied to stored SBOMs\\
\checkbox & Activities and Tools mapping of REQUIRED and PREFERRED activities, with a roadmap and demonstrations\\
\checkbox & CNAPP: composition analysis, SAST/DAST/IAST, cloud posture and entitlement management, IaC scanning, workload protection, cloud detection and response\\
\midrule \multicolumn{2}{@{}l}{\sffamily\bfseries Process \pg{cATO criteria}{15--16}}\\ \midrule
\checkbox & IaC and configuration as code to prevent drift\\
\checkbox & Each control gate described (what opens and closes it, alerts, responses) and demonstrated; each guardrail and its out-of-tolerance process\\
\midrule \multicolumn{2}{@{}l}{\sffamily\bfseries People \pg{cATO criteria}{16--17}}\\ \midrule
\checkbox & Role-based training records; org chart; separation of duties and least privilege\\
\checkbox & Tabletop exercises with after-action reports (incidents, gate triggers, alerts, privilege requests)\\
\checkbox & Team trained on the reference design, tools, gates and tolerances, adjudication, root cause analysis, POA\&Ms and dashboards\\
\checkbox & Insider threat working group; onboarding and offboarding for everyone regardless of rank\\
\end{longtable}

\begin{boxteam}{cATO}
A small model team rarely seeks its own cATO; it builds on a platform that has (or is pursuing) one, and inherits most
of the platform controls. What it owns is in its repository: control gates that are documented and demonstrable (the
pipeline jobs and their failure messages), an SBOM and an archive of them per baseline, findings tracked to closure in
the backlog, and training records. Ask the platform team for its reference design, its risk tolerances and how a
pipeline failure is routed (Primer: \emph{Security: protect the environment, engineer it into the system}).
\end{boxteam}

% =====================================================================================================================
\chapter{What it means for a small MBSE team}\label{ch:team}

This chapter is this guide's reading, not the Department's. It assumes a team like the MBSE-DEVSECOPS Primer's: a few
systems engineers keeping a SysML v2 text model in GitLab Dedicated, runners on Kubernetes (AKS) in an Azure
subscription, the environment defined in Terraform, work arriving as requests from another system.

\begin{longtable}{@{}L{0.33\linewidth}L{0.38\linewidth}L{0.23\linewidth}@{}}
\caption{From obligation to practice.}\label{tab:team}\\
\toprule \sffamily\bfseries The documents ask & \sffamily\bfseries A model team does & \sffamily\bfseries Primer chapter\\ \midrule\endfirsthead
\toprule \sffamily\bfseries The documents ask & \sffamily\bfseries A model team does & \sffamily\bfseries Primer chapter\\ \midrule\endhead
\bottomrule\endfoot
Testable requirements traceable to mission functions \pg{DoWI 8430.01}{4--5} & Requirements as model elements with satisfy and verify links, checked on every merge request & Requirements a pipeline can check\\
Operate a secure CI/CD pipeline \pg{DoWI 8430.01}{9} & Every model change through a merge request and a pipeline; protected branches & The pipeline, job by job\\
Testing is continuous, not a final gate \pg{DoWI 8430.01}{14} & Grammar, lint, trace and verification-matrix jobs on every change; results kept as evidence & CI/CD from zero\\
Configuration management discipline \pg{Guidebook}{15} & Baselines as tags, audited by checksum; the change history in git & Maintenance through the life cycle\\
SBOM and provenance \pg{DoWI 8430.01}{18} & Record the toolchain (versions and checksums) with each baseline; generate an SBOM for any software produced & Security: protect the environment\\
Value stream metrics; DORA inside the team \pg{DoWI 8430.01}{15--16} & Lead time from request to merged change and to baseline; deployment frequency of baselines; failure rate of pipelines & Measuring with GitLab analytics\\
Engage users; demonstrate \pg{DoWI 8430.01}{10} & Sprint reviews that demonstrate the model, not slides about it & The rhythm\\
AI code is unverified input \pg{DoWI 8430.01}{23} & Review every AI-drafted change like any other; label it; keep the assistant on an approved system & A local AI assistant on the laptop\\
Use an existing platform \pg{cATO guide}{15} & Use the enterprise GitLab and runners; add jobs, not infrastructure & Setting up GitLab; Runners\\
Infrastructure as Code \pg{Strategy}{11} & Keep the project's own settings (labels, boards, schedules) as code; leave the cluster's Terraform to the platform team unless you own it & Administration\\
\end{longtable}

\begin{boxnote}{What not to over-read}
None of the documents requires a particular tool, a SysML language or a model repository. The guidebook says common
tools are ``a class of tooling'' and not the same product everywhere \pg{Fundamentals}{20}. The DORA metrics and the
twelve-week scorecard are team diagnostics; reporting upward should be in mission terms \pg{DoWI 8430.01}{16},
\pg{State report}{41}.
\end{boxnote}

% =====================================================================================================================
\chapter{What to read in full, and when}\label{ch:read}

\begin{longtable}{@{}L{0.22\linewidth}L{0.36\linewidth}L{0.36\linewidth}@{}}
\caption{Reading plan.}\label{tab:read}\\
\toprule \sffamily\bfseries Who & \sffamily\bfseries Read in full & \sffamily\bfseries When\\ \midrule\endfirsthead
\toprule \sffamily\bfseries Who & \sffamily\bfseries Read in full & \sffamily\bfseries When\\ \midrule\endhead
\bottomrule\endfoot
Everyone on the team & DevSecOps Fundamentals (44 pp.); DoWI 8430.01 section 3 (pp.~13--29) & In the first two weeks\\
Systems engineering lead & DoWI 8430.01 in full; Activities \& Tools Guidebook Plan, Develop and Test tables & Before the first planning meeting; again before a review\\
Whoever owns the pipeline & Guidebook and workbook in full; the reference design your platform implements & Before writing the pipeline; when the platform changes\\
Security and the ISSM & cATO memo, evaluation criteria, implementation guide; DoWI 8430.01 section 3.5 & Before the authorization conversation\\
Program manager & Strategy; Implementation Plan; State of DevSecOps sections 5 and 8; DoWI 8430.01 sections 3.3 and 3.8 & Before setting measures and contracts\\
Anyone choosing a tool or library & OSS memo; DoWI 8430.01 section 3.5.i & Before adopting it\\
Anyone using an AI assistant & DoWI 8430.01 section 3.6 & Before the first use\\
For background only & Playbook, Pathway, SWFT memo and RFI summary, AWS draft & When the topic comes up\\
\end{longtable}

% =====================================================================================================================
\appendix
\chapter{Glossary of acronyms}\label{app:glossary}
\begin{multicols}{2}\small
\begin{description}[style=nextline,leftmargin=0pt,itemsep=1pt,font=\sffamily\bfseries]
\item[ACD] Active Cyber Defense
\item[AO / AODR] Authorizing Official / AO Designated Representative
\item[ATO / ATC] Authorization to Operate / to Connect
\item[cATO] Continuous Authorization to Operate
\item[CI/CD] Continuous Integration / Continuous Delivery (or Deployment)
\item[CISO] Chief Information Security Officer
\item[CM / CMDB] Configuration Management / CM Database
\item[CNAP] Cloud Native Access Point
\item[CNAPP] Cloud Native Application Protection Platform
\item[CNCF] Cloud Native Computing Foundation
\item[CONMON] Continuous Monitoring
\item[CSSP] Cybersecurity Service Provider
\item[CTI] Controlled Technical Information
\item[CUI] Controlled Unclassified Information
\item[DAST / SAST / IAST] Dynamic / Static / Interactive Application Security Testing
\item[DCWF] DoD (DoW) Cyber Workforce Framework
\item[DITPR] DoD IT Portfolio Repository
\item[DORA] DevOps Research and Assessment
\item[DSOP] DevSecOps Platform
\item[DT\&E / OT\&E] Developmental / Operational Test and Evaluation
\item[E2P] Exception to Policy
\item[IaC / CaC] Infrastructure / Configuration as Code
\item[IL] Impact Level (cloud security requirements guide)
\item[ISA / MOU] Interconnection Security Agreement / Memorandum of Understanding
\item[ISCM] Information Security Continuous Monitoring
\item[JWCC] Joint Warfighting Cloud Capability
\item[MOSA] Modular Open Systems Approach
\item[MVP / MVCR] Minimum Viable Product / Capability Release
\item[NSS] National Security System
\item[OSS] Open Source Software
\item[POA\&M] Plan of Action and Milestones
\item[RAR / SAR / SAP / SSP] Risk Assessment Report / Security Assessment Report / Plan / System Security Plan
\item[RMF / RMF KS] Risk Management Framework / its Knowledge Service
\item[SBOM / SCA] Software Bill of Materials / Software Composition Analysis
\item[SCSS] Sidecar Container Security Stack
\item[SIEM / SOAR] Security Information and Event Management / Security Orchestration, Automation and Response
\item[SSDF] Secure Software Development Framework (NIST SP 800-218); groups PO, PS, PW, RV
\item[SSG] Software Modernization Senior Steering Group
\item[SSSC] Secure Software Supply Chain
\item[SWFT] Software Fast Track
\item[SWP] Software Acquisition Pathway (DoDI 5000.87)
\end{description}
\end{multicols}

\chapter{Sources}\label{app:sources}
The library page \url{https://dodcio.defense.gov/library/} as captured on 3 October 2026 (the live site did not answer
automated requests when this guide was compiled). Pages are PDF pages. Checksums are SHA-256 of the copies read.

{\small
\begin{longtable}{@{}L{0.04\linewidth}L{0.90\linewidth}@{}}
\toprule \sffamily\bfseries \# & \sffamily\bfseries Document, address, version, pages, SHA-256\\ \midrule\endfirsthead
\toprule \sffamily\bfseries \# & \sffamily\bfseries Document, address, version, pages, SHA-256\\ \midrule\endhead
\bottomrule\endfoot
1 & \textbf{DoW Software Modernization Strategy} --- \url{https://dodcio.defense.gov/Portals/0/Documents/Library/SoftwareModStrat.pdf} --- v1.0 Nov 2021, memo 1 Feb 2022 --- 14 pp. --- \texttt{\scriptsize 0276c47b7fb058e660bdacce2ca621e93a8cea78e290a7b05f9f581b8a02b53a}\\
2 & \textbf{Software Modernization Implementation Plan FY25--26} --- \url{https://dodcio.defense.gov/Portals/0/Documents/Library/SW-Mod-I-Plan25-26.pdf} --- cleared 30 Apr 2025 --- 27 pp. --- \texttt{\scriptsize ba350c72245b8efac4a0e4a5655f4042e9a53b6e8a1fa44dbe3170bae0122601}\\
3 & \textbf{DoWI 8430.01, Accelerated Mission Software} --- \url{https://www.esd.whs.mil/Portals/54/Documents/DD/issuances/dodi/843001p.pdf} --- effective 8 Sep 2026 --- 37 pp. --- \texttt{\scriptsize b35f0dcc2fae721994413c702671e9b754ed6911daeecfcfec0d83f0a4df465a}\\
4 & \textbf{Memo: Accelerating Secure Software} --- \url{https://dodcio.defense.gov/Portals/0/Documents/Library/Memo-AcceleratingSecureSoftware.pdf} --- 24 Apr 2025 --- 1 p. --- \texttt{\scriptsize 2010080899dceab07f1dd9caa48ed5cfba6d921eccc34cd5a7b790efdabdea13}\\
5 & \textbf{Software Fast Track RFI Combined Summary} --- \url{https://dodcio.defense.gov/Portals/0/Documents/Library/SWFT-RFI-Combined-Summary.pdf} --- cleared 3 Dec 2025 --- 31 pp. --- \texttt{\scriptsize a2a2ad2a36b9e336775e567e3a6c60ba6745f8ef204d5bf9ab0066dfc06825b1}\\
6 & \textbf{DoW Enterprise DevSecOps Fundamentals} --- \url{https://dodcio.defense.gov/Portals/0/Documents/Library/DoD%20Enterprise%20DevSecOps%20Fundamentals%20v2.5.pdf} --- v2.5, Oct 2024 --- 44 pp. --- \texttt{\scriptsize dbef8ab70c33aade3d6495970caad09f46f750482625e8ad22400c7b1eddb22b}\\
7 & \textbf{DoW Enterprise DevSecOps Activities \& Tools Guidebook} --- \url{https://dodcio.defense.gov/Portals/0/Documents/Library/DevSecOpsActivitesToolsGuidebook.pdf} --- v2.5, Apr 2025 --- 65 pp. --- \texttt{\scriptsize 350432bf6bc046a8a0ab5e5ece352386dabcb8a7921c54fcacd9b49897224143}\\
8 & \textbf{DevSecOps Fundamentals Guidebook: Activities \& Tools (xlsx)} --- \url{https://dodcio.defense.gov/Portals/0/Documents/Library/DevSecOpsActivitesToolsGuidebookTables.xlsx} --- companion of v2.5 --- 14 tabs --- \texttt{\scriptsize 60c330c7659d2fa242deba078722c5c8242883c9c02cc0f147506c2459b623a9}\\
9 & \textbf{DevSecOps Playbook} --- \url{https://dodcio.defense.gov/Portals/0/Documents/Library/DevSecOps%20Playbook_DoD-CIO_20211019.pdf} --- v2.1, Sep 2021 --- 14 pp. --- \texttt{\scriptsize 809a6fbb462aeeec20b4db087b3feddf084dcd9ff855842e7b18e990ce882554}\\
10 & \textbf{The State of DevSecOps Report} --- \url{https://dodcio.defense.gov/Portals/0/Documents/Library/DevSecOpsStateOf.pdf} --- Mar 2025 --- 47 pp. --- \texttt{\scriptsize 66a647df3b1b1bf837e430e5aff1fb4b0a0370979ff55ec4d30f0c9a11f3874e}\\
11 & \textbf{Reference Design: CNCF Kubernetes} --- \url{https://dodcio.defense.gov/Portals/0/Documents/Library/DoD%20Enterprise%20DevSecOps%20Reference%20Design%20-%20CNCF%20Kubernetes%20w-DD1910_cleared_20211022.pdf} --- v2.1, Sep 2021 --- 32 pp. --- \texttt{\scriptsize 9742f6a5d6ca6b32ca30c6857a61e3d35d009141df2d5e6891a77eaf7164a117}\\
12 & \textbf{Reference Design: CNCF Multi-Cluster Kubernetes} --- \url{https://dodcio.defense.gov/Portals/0/Documents/Library/DoDReferenceDesign-CNCFMulti-ClusterKubernetes.pdf} --- v1.0, Jul 2022 --- 28 pp. --- \texttt{\scriptsize addad98c2d19c620f4ef0260687595b731e0d11fb9eef93f9fe3e460ef7ffef1}\\
13 & \textbf{DRAFT: DOW AWS Managed Service} --- \url{https://dodcio.defense.gov/Portals/0/Documents/Library/DoD%20Enterprise%20DevSecOps%20Reference%20Design%20-%20AWS%20Managed%20Services_DoD-CIO_20211019.pdf} --- v0.2, Sep 2021 --- 35 pp. --- \texttt{\scriptsize 0041626903b05b2cd8fbb3e8aa323fe5de4624d29148d17425498bc0aeabef9e}\\
14 & \textbf{DRAFT: DOW Cloud IaC + GitHub + Azure} --- \url{https://dodcio.defense.gov/Portals/0/Documents/Library/DoDRefDesignCloudGithub.pdf} --- v0.01, Sep 2022 --- 70 pp. --- \texttt{\scriptsize 2e2e896d30588754cc0be268ac896b0e5c8b35df6c93a843a990ff633838fc0d}\\
15 & \textbf{Pathway to Reference Design} --- \url{https://dodcio.defense.gov/Portals/0/Documents/Library/DoD%20Enterprise%20DevSecOps-Pathway%20to%20a%20Reference%20Design_DoD-CIO_20211018.pdf} --- Sep 2021 --- 3 pp. --- \texttt{\scriptsize 0e4b2f85eb0d8eaeda3098f7aa4c2ba17298fcb1b36792aa19952e0934ca733b}\\
16 & \textbf{Software Development and Open Source Software Memorandum} --- \url{https://dodcio.defense.gov/Portals/0/Documents/Library/SoftwareDev-OpenSource.pdf} --- 24 Jan 2022 --- 8 pp. --- \texttt{\scriptsize 8113f7c6321f5adf549de678bfc8bfb60fd784b03cbcf73d34f88ae7e6345cff}\\
17 & \textbf{Memo: Continuous Authorization to Operate (cATO)} --- \url{https://dodcio.defense.gov/Portals/0/Documents/Library/20220204-cATO-memo-Signed-Cleared.pdf} --- 2 Feb 2022 --- 3 pp. --- \texttt{\scriptsize c366edb4bd0174574c12dae25362a40cd7428f871dc2a0676cf76d2da6580c90}\\
18 & \textbf{DevSecOps cATO Evaluation Criteria} --- \url{https://dodcio.defense.gov/Portals/0/Documents/Library/cATO-EvaluationCriteria.pdf} --- 29 May 2024 --- 19 pp. --- \texttt{\scriptsize afbe0a13a7f63c8741d036d373508c93ae1c9df1abb37ddf74714f1b43052608}\\
19 & \textbf{DevSecOps Continuous Authorization Implementation Guide} --- \url{https://dodcio.defense.gov/Portals/0/Documents/Library/DoDCIO-ContinuousAuthorizationImplementationGuide.pdf} --- v1.0, Mar 2024 --- 26 pp. --- \texttt{\scriptsize 4218c2c20efd5002ea107e1f74d9e7777009fe9efe66f04a5bbaba142918712a}\\
\end{longtable}}

\end{document}
